\scanpage{165}
\mspage{79}
\[
(267)\qquad X=\mathbb C^I,\qquad \Theta=\operatorname{div}\Bigl(\prod_{i\in I}z_i\Bigr)\deleted{\unsure{\mathrm{red}}}
\]
où $I$ est un ens. fini de cardinal $n$, et où $(z_i)_{i\in I}$ sont les fonctions coordonnées.

On s'attend à ce que les ``topos élémentaires'' qui correspondent à cette situation soient les produits (indexés par $I$) des topos de l'un quelc. des quatre types envisagés dans (264) -- et que les morphismes canoniques entre ces topos doivent \deleted{\unsure{correspondre}}\inserted{se déduire} des morphismes inclus dans le diagramme (265). Cela nous \deleted{suggère}\inserted{conduit} à considérer les partitions possibles de $I$ en quatre paquets, correspondant aux quatre types en question. On écrira
\margin{``partition'' \unsure{sens généralisé} -- il peut y avoir des ``classes'' vides}
\[
(268)\quad\unsure{\mathrm{a}}\qquad
\left\{\begin{array}{l}
I=\underbrace{I'}_{\substack{\text{correspond}\\\text{aux facteurs}\\\text{de type}\\\text{``global''}}}\amalg
\underbrace{I''}_{\substack{\text{correspond}\\\text{aux facteurs}\\\text{de type local}}}\\[3em]
I'=\underbrace{I'_0}_{\substack{\text{correspond}\\\text{aux facteurs }\mathbb P\\\text{(de dim }0\text{)}}}\amalg
\underbrace{I'_1}_{\substack{\text{correspond}\\\text{aux facteurs }\mathbb C^*\\\text{(de dim }1\text{)}}}\\[3em]
I''=\underbrace{I''_0}_{\text{facteurs }\mathbb D}\amalg
\underbrace{I''{}^*}_{\text{facteurs }\mathbb D^*\text{ (épointés)}}
\end{array}\right.
\]
\margin{164}
